\section{Literatür verisiyle yazılım uygulamaları}
\label{sec:spss-b11}

\textbf{Amaç ve veri.} \texttt{b11.csv} ile iki ayrı soru sorulur:
F ve M kodlu grupların yıl sonu ortalamaları farklı mı; aynı öğrencinin
\texttt{g3} ve \texttt{g1} notları arasında ortalama fark var mı
\citep{cortez2008data}?
İlk soru için Welch, ikinci soru için eşleştirilmiş $t$ testi kullanılır.
İki test de iki yönlüdür; her biri için anlamlılık düzeyi 0,05 ve
güven düzeyi yüzde 95'tir. Fark yönleri sırasıyla F eksi M ve
\texttt{g3-g1} olarak sabitlenmiştir.

\textbf{Ortak veri.} Doğrulanan B11 uygulamasında üç yazılım da
\nolinkurl{companion/spss/csv/b11.csv} dosyasını kullanmıştır.
Dosya 395 satır ve 10 sütun içerir; kayıtların tamamı analiz edilir.
\texttt{sex} için 1=F ve 2=M kaynak kategorileridir. B06'nın seçim
göstergeleri ve ağırlığı kullanılmaz. Bu karşılaştırma Excel içe
aktarmanın doğrulandığı anlamına gelmez. Kodları \texttt{companion}
klasörünü içeren paket kökünden çalıştırın; veri sözlüğü ve hazırlık
için bkz. sayfa~\pageref{sec:spss-guide}.

\subsection{IBM SPSS uygulaması}
\textbf{Başlamadan önce.} B10'dan kalan tek satırlık özeti değil,
aşağıdaki veri açma bloğuyla ham B11 verisini kullanın. Menü yoluna
geçmeden önce bu bloğu \texttt{EXECUTE} satırı dahil çalıştırın.
Filtre, ağırlık ve dosya bölme kapalı olmalıdır. Dosya açma hatasında
önce \texttt{/FILE} yolunu düzeltip veri açma bloğunu yeniden çalıştırın.
B11 kodu kayıtları özet satırına indirmez; yalnız \texttt{fark}
değişkenini ekler.

\begin{spssadimlar}
\spssadim{\emph{Analyze $\to$ Descriptive Statistics $\to$ Explore} yolunda \texttt{g3} değişkenini \emph{Dependent List}, \texttt{sex} değişkenini \emph{Factor List} alanına yerleştirin; grup kutu grafiklerini inceleyin.}
\spssadim{\emph{Analyze $\to$ Compare Means $\to$ Independent-Samples T Test} yolunu açın. \texttt{g3} test değişkeni, \texttt{sex} gruplama değişkeni olsun.}
\spssadim{\emph{Define Groups} içinde grup 1 için 1, grup 2 için 2 yazın. \emph{Continue}; yüzde 95 güven düzeyini kontrol edip \emph{OK} seçin.}
\spssadim{Önceden seçilen Welch yaklaşımı için \emph{Equal variances not assumed} satırını ve iki yönlü anlamlılık sütununu okuyun. Farkın yönü F eksi M'dir. Levene sonucuna göre sonradan satır değiştirmeyin.}
\spssadim{Eşleştirilmiş analiz öncesi \emph{Transform $\to$ Compute Variable} ile \texttt{fark=g3-g1} hesaplayın; Explore ile bu farkın histogram ve kutu grafiğini inceleyin.}
\spssadim{\emph{Analyze $\to$ Compare Means $\to$ Paired-Samples T Test} yolunda ilk değişken \texttt{g3}, ikinci değişken \texttt{g1} olacak biçimde tek çift oluşturun.}
\spssadim{\emph{Options} içinde yüzde 95 güven düzeyini seçip \emph{Continue $\to$ OK} ile çalıştırın. Fark \texttt{g3-g1} yönündedir. Yeni sürümlerde \emph{Compare Means} yerine \emph{Compare Means and Proportions} menüsü kullanılabilir.}
\end{spssadimlar}

{\raggedright
\noindent\textbf{Komutların görevi.} \texttt{GROUPS=sex(1 2)} bağımsız grupları ve fark yönünü tanımlar. \texttt{PAIRS=g3 WITH g1} aynı öğrencinin iki notunu eşleştirir. Aradaki \texttt{COMPUTE} ve \texttt{EXAMINE}, fark dağılımını kontrol eder.\par}

\par\noindent\begin{minipage}{\linewidth}
\textbf{Uygulama kodu:} \texttt{b11}. Doğrulanan veri dosyası
\texttt{companion/spss} altındaki \texttt{csv/b11.csv} dosyasıdır.

\begin{lstlisting}[style=prismSPSS]
SET DECIMAL=DOT.
GET DATA /TYPE=TXT
 /FILE='companion/spss/csv/b11.csv'
 /ENCODING='UTF8'
 /ARRANGEMENT=DELIMITED /DELCASE=LINE
 /DELIMITERS="," /QUALIFIER='"'
 /FIRSTCASE=2
 /VARIABLES=id F8.0 school F8.0 sex F8.0 age F8.0
 studytime F8.0 absences F8.0 g1 F8.0 g2 F8.0
 g3 F8.0 pass10 F8.0.
FILTER OFF.
WEIGHT OFF.
SPLIT FILE OFF.
VARIABLE LABELS
 sex 'Kaynak verideki cinsiyet kategorisi'
 g1 'Ilk donem matematik notu'
 g3 'Yil sonu matematik notu'.
VALUE LABELS sex 1 'F' 2 'M'.
VARIABLE LEVEL
 id school sex pass10 (NOMINAL)
 studytime (ORDINAL)
 age absences g1 g2 g3 (SCALE).
FORMATS g1 g3 (F14.9).
EXECUTE.
\end{lstlisting}
\end{minipage}\par

\par\noindent\begin{minipage}{\linewidth}
\textbf{İki analiz.} Veri açma bloğundan sonra çalıştırın.

\begin{lstlisting}[style=prismSPSS]
EXAMINE VARIABLES=g3 BY sex
 /PLOT=BOXPLOT /STATISTICS=DESCRIPTIVES
 /CINTERVAL=95.
T-TEST GROUPS=sex(1 2) /VARIABLES=g3
 /CRITERIA=CI(.95).
COMPUTE fark=g3-g1.
VARIABLE LABELS fark 'Yil sonu eksi ilk donem notu'.
VARIABLE LEVEL fark (SCALE).
FORMATS fark (F14.9).
EXECUTE.
EXAMINE VARIABLES=fark
 /PLOT=BOXPLOT HISTOGRAM
 /STATISTICS=DESCRIPTIVES /CINTERVAL=95.
T-TEST PAIRS=g3 WITH g1 (PAIRED)
 /CRITERIA=CI(.95).
\end{lstlisting}

\end{minipage}\par

\begin{outputguidebox}
\textbf{Paylaşılan SPSS tablolarıyla doğrulanan değerler.}
\emph{Group Statistics} tablosunda F için $n=208$, ortalama
9,9663461538; M için $n=187$, ortalama 10,914438503 görülür.
\emph{Equal variances not assumed} satırında $t=-2{,}065$,
serbestlik derecesi 390,569 ve iki yönlü $p=0{,}040$ vardır.
F eksi M farkı $-0{,}9480923488$ puandır. Eşit varyans varsayılan
satırın aynı yuvarlanmış $p$ değerini vermesi, iki yöntemi aynı yapmaz.

\emph{Paired Samples Test} tablosunda \texttt{g3-g1} farkı
$-0{,}4936708861$, $t(394)=-3{,}552$ ve iki yönlü $p<0{,}001$
görülür. Bu ifade $p=0$ demek değildir. \emph{Paired Samples
Correlations} tablosundaki 0,801 korelasyon ve onun anlamlılığı,
ortalama farkının testi yerine okunmaz. Her öğrencinin iki notu
bir çifttir; 790 bağımsız öğrenci yoktur.
\end{outputguidebox}

\textbf{Grafik kontrolünün kapsamı.} Paylaşılan SPSS fark histogramında
395 kayıt, ortalama $-0{,}493670886$ ve standart sapma 2,762481677
görülür. Farklar sıfır çevresinde yoğunlaşırken negatif tarafta uzun
bir kuyruk vardır. Farkın kutu grafiğinde medyan sıfır civarındadır;
negatif tarafta çok sayıda aykırı değer işaretlenmiştir. Bu görseller
simetrik normal dağılım görünümünü desteklemez; normallik doğrulanmış
sayılmaz. İşaretlenen kayıtlar yalnız aykırı etiketi nedeniyle silinmez.

Sonradan paylaşılan SPSS F--M \texttt{g3} kutu grafiği de incelenmiştir.
F grubunda medyan 10, kutu sınırları 8 ve 13; M grubunda medyan 11,
kutu sınırları 9 ve 14'tür. Bıyıklar sırasıyla 4--19 ve 5--20 puana
uzanır. M grubunun merkezi daha yüksek görünse de dağılımlar belirgin
biçimde örtüşür; kutudaki çizgi ortalamayı değil medyanı gösterir.
Her iki grupta sıfır notlar aykırı olarak işaretlenmiştir. Üst üste
binen simgelerden kayıt sayısı okunmaz; proje CSV'sinde F grubunda
23, M grubunda 15 sıfır not vardır ve bu 38 kayıt analizde tutulur.
Bu görünüm tek başına normalliği, varyans eşitliğini veya gözlemlerin
bağımsızlığını doğrulamaz; önceden seçilen Welch testi değiştirilmez.


\par\noindent\begin{minipage}{\linewidth}
\subsection{R uygulaması}
\label{sec:r-gercek-b11}

\texttt{var.equal=FALSE} Welch testini seçer. Eşleştirilmiş testte
ilk değişken G3, ikinci değişken G1 olduğundan fark G3--G1 olur.
Aşağıdaki R bloklarını sırayla çalıştırın.

\begin{lstlisting}[style=prismR]
veri <- read.csv("companion/spss/csv/b11.csv")
stopifnot(nrow(veri) == 395, ncol(veri) == 10,
  !anyNA(veri), !anyDuplicated(veri$id),
  all(veri$sex %in% c(1, 2)),
  all(veri$pass10 == as.integer(veri$g3 >= 10)))
grup_F <- veri$g3[veri$sex == 1]
grup_M <- veri$g3[veri$sex == 2]
fark <- veri$g3-veri$g1
welch <- t.test(grup_F, grup_M, var.equal=FALSE,
  alternative="two.sided", conf.level=.95)
eslesmis <- t.test(veri$g3, veri$g1, paired=TRUE,
  alternative="two.sided", conf.level=.95)
print(welch); print(eslesmis)
\end{lstlisting}

\textbf{Çıktıyı okuma.} Paylaşılan R çıktısında Welch için iki yönlü
$p=0{,}039577$, eşleştirilmiş test için $p=0{,}000429$ bulunmuştur.
İki testin güven aralıkları da ilgili ortalama farkına aittir.
\end{minipage}\par

\par\noindent\begin{minipage}{\linewidth}
\textbf{Ayrıntılı özetler.} Aşağıdaki kod grup ve ölçüm özetlerini,
iki testin sonuçlarını ve eşleştirilmiş etkiyi yazdırır.

\begin{lstlisting}[style=prismR]
betimsel <- function(degerler) {
  c(n=length(degerler), ort=mean(degerler),
    ss=sd(degerler), sh=sd(degerler)/sqrt(length(degerler)))
}
test_ozeti <- function(sonuc, ortalama_fark) {
  c(t=unname(sonuc$statistic), df=unname(sonuc$parameter),
    p_iki_yonlu=sonuc$p.value, fark=ortalama_fark,
    alt=sonuc$conf.int[1], ust=sonuc$conf.int[2])
}
print(rbind(F=betimsel(grup_F), M=betimsel(grup_M)),
      digits=15)
print(test_ozeti(welch, mean(grup_F)-mean(grup_M)),
      digits=15)
print(rbind(G3=betimsel(veri$g3), G1=betimsel(veri$g1)),
      digits=15)
print(test_ozeti(eslesmis, mean(fark)), digits=15)
print(c(betimsel(fark), dz=mean(fark)/sd(fark)), digits=15)
print(sum(veri$g3 == 0))
\end{lstlisting}

\textbf{Çıktıyı okuma.} R çıktısında 395 satır, 10 sütun ve sıfır
eksik değer vardır; 38 sıfır not analizde tutulmuştur.
$d_z=-0{,}178706$, ortalama farkın \emph{farkların standart sapmasına}
bölünmesiyle hesaplanır; iki grubun ortak standart sapması kullanılmaz.
\end{minipage}\par

\par\noindent\begin{minipage}{\linewidth}
\textbf{Grafikleri üretme.} Bu R grafiklerinin çıktısı paylaşılmamıştır;
görsel inceleme için kod aşağıda ayrıca verilmiştir.

\begin{lstlisting}[style=prismR]
par(mfrow=c(1, 3))
boxplot(g3 ~ sex, data=veri, names=c("F", "M"))
boxplot(fark, main="G3-G1")
hist(fark, main="G3-G1", xlab="Difference")
par(mfrow=c(1, 1))
\end{lstlisting}

\end{minipage}\par

\par\noindent\begin{minipage}{\linewidth}
\subsection{Python uygulaması}
\label{sec:python-b11}

\texttt{equal\_var=False} Welch testidir; \texttt{ttest\_rel} aynı
öğrencinin iki notunu eşleştirir. Python bloklarını sırayla çalıştırın;
iki testin aralıkları ayrı okunur.

\begin{lstlisting}[style=prismPythonApp]
import pandas as pd
import numpy as np
from scipy import stats

veri = pd.read_csv("companion/spss/csv/b11.csv")
assert veri.shape == (395, 10)
assert not veri.isna().any().any()
assert veri.id.is_unique
assert veri.sex.isin([1, 2]).all()
assert veri.pass10.eq((veri.g3 >= 10).astype(int)).all()
grup_F = veri.loc[veri.sex == 1, "g3"]
grup_M = veri.loc[veri.sex == 2, "g3"]
fark = veri.g3-veri.g1
welch = stats.ttest_ind(grup_F, grup_M,
    equal_var=False, alternative="two-sided")
eslesmis = stats.ttest_rel(veri.g3, veri.g1,
                          alternative="two-sided")
print(welch); print(welch.confidence_interval(.95))
print(eslesmis); print(eslesmis.confidence_interval(.95))
\end{lstlisting}

\textbf{Çıktıyı okuma.} Paylaşılan Python çıktısında fark yönleri
F--M ve G3--G1 olarak R ve SPSS ile aynıdır. Serbestlik derecesi
Welch testinde kesirli, eşleştirilmiş testte 394'tür.
\end{minipage}\par

\par\noindent\begin{minipage}{\linewidth}
\textbf{Ayrıntılı özetler.} \texttt{ddof=1}, örneklem standart sapmasını
seçer. Aşağıdaki özetlerin sayısal içeriği yüklenen Python çıktısıyla
aynıdır; burada yazdırma düzeni kısaltılmıştır.

\begin{lstlisting}[style=prismPythonApp]
def betimsel(degerler):
    return dict(n=len(degerler), ort=degerler.mean(),
        ss=degerler.std(ddof=1),
        sh=degerler.std(ddof=1)/np.sqrt(len(degerler)))

def test_ozeti(sonuc, ortalama_fark):
    aralik = sonuc.confidence_interval(.95)
    return dict(t=sonuc.statistic, df=sonuc.df,
        p_iki_yonlu=sonuc.pvalue, fark=ortalama_fark,
        alt=aralik.low, ust=aralik.high)

ozetler = {
    "F": betimsel(grup_F), "M": betimsel(grup_M),
    "Welch": test_ozeti(welch, grup_F.mean()-grup_M.mean()),
    "G3": betimsel(veri.g3), "G1": betimsel(veri.g1),
    "Eslesmis": test_ozeti(eslesmis, fark.mean()),
    "Fark": dict(betimsel(fark),
                 dz=fark.mean()/fark.std(ddof=1))}
for baslik, ozet in ozetler.items():
    print(baslik)
    for ad, deger in ozet.items():
        print(f"{ad}: {deger:.13f}")
print(int((veri.g3 == 0).sum()))
\end{lstlisting}

\textbf{Çıktıyı okuma.} Python çıktısı da 395 satır, 10 sütun, sıfır
eksik değer ve 38 sıfır not gösterir. R ve Python ham kayıtları korur.
\end{minipage}\par

\par\noindent\begin{minipage}{\linewidth}
\textbf{Grafikleri üretme.} Bu Python grafiklerinin çıktısı
paylaşılmamıştır; sayısal uyum grafiklerin doğrulandığı anlamına gelmez.

\begin{lstlisting}[style=prismPythonApp]
import matplotlib.pyplot as plt

sekil, eksenler = plt.subplots(1, 3, figsize=(12, 4))
eksenler[0].boxplot([grup_F, grup_M])
eksenler[0].set_xticks([1, 2], ["F", "M"])
eksenler[0].set_title("G3")
eksenler[1].boxplot(fark); eksenler[1].set_title("G3-G1")
eksenler[2].hist(fark, bins=15)
eksenler[2].set_title("G3-G1")
plt.tight_layout(); plt.show()
\end{lstlisting}

\end{minipage}\par

\Needspace{15\baselineskip}
\subsection{Sonuçların karşılaştırılması ve ortak rapor}
13 Eylül 2026'da paylaşılan SPSS tabloları, R ve Python metin
çıktılarıyla karşılaştırılmıştır. Beş betimsel özetin dörder değeri,
iki testin altışar değeri ve $d_z$ olmak üzere 33 sayısal özetin
tamamı R ve Python arasında $10^{-12}$ mutlak tolerans içinde
uyumludur. Python'un 13 ondalık basamaklı özetleri proje CSV'sinden
yapılan hesaplarla da eşleşir. SPSS'in gösterdiği basamaklar aynı
sonuçlarla uyumludur; eşleştirilmiş testin $p<0{,}001$ gösterimi,
R/Python'daki yaklaşık 0,000429 değerini içerir.
Görsel kontrol paylaşılan SPSS F--M grup kutu grafiğini, fark
histogramını ve fark kutu grafiğini kapsar. R/Python grafiklerinin
çıktıları paylaşılmadığından bunlar görsel doğrulama kapsamında
değildir. Aşağıdaki tablolar çıktıların kitap
düzenine uyarlanmış özetidir; ekran görüntüsü değildir.

\begin{table}[H]
\centering
\caption{B11'in doğrulanmış grup, ölçüm ve fark özetleri}
\label{tab:b11-dogrulanmis-betimsel}
\begin{tabular}{@{}lrrrr@{}}
\toprule
Değişken/grup & $n$ & Ortalama & Standart sapma & Standart hata \\
\midrule
G3 (F) & 208 & 9,966346 & 4,622338 & 0,320501 \\
G3 (M) & 187 & 10,914439 & 4,495297 & 0,328729 \\
\midrule
G3 (tümü) & 395 & 10,415190 & 4,581443 & 0,230517 \\
G1 (tümü) & 395 & 10,908861 & 3,319195 & 0,167007 \\
G3--G1 & 395 & $-0{,}493671$ & 2,762482 & 0,138996 \\
\bottomrule
\end{tabular}
\par\smallskip
{\footnotesize Ondalıklı değerler altı basamağa yuvarlanmıştır.
İlk iki satır aynı 395 öğrencinin F ve M gruplarını, son üç satır
bu öğrencilerin iki notunu ve kişi içi farkını özetler; satırlar
ayrı örneklemler olarak toplanmaz. Standart hata $s/\sqrt{n}$'dir.\par}
\end{table}

\begin{table}[H]
\centering
\caption{B11'de Welch ve eşleştirilmiş test sonuçları}
\label{tab:b11-dogrulanmis-test}
\begin{tabular}{@{}lrr@{}}
\toprule
Ölçü & Welch (F--M) & Eşleştirilmiş (G3--G1) \\
\midrule
Ortalama fark & $-0{,}948092$ & $-0{,}493671$ \\
$t$ istatistiği & $-2{,}065057$ & $-3{,}551703$ \\
Serbestlik derecesi & 390,569356 & 394 \\
İki yönlü $p$ değeri & 0,039577 & 0,000429 \\
\%95 aralık: alt sınır & $-1{,}850732$ & $-0{,}766937$ \\
\%95 aralık: üst sınır & $-0{,}045452$ & $-0{,}220405$ \\
\bottomrule
\end{tabular}
\par\smallskip
{\footnotesize Aralıklar ilgili ortalama farkları içindir; tek bir
grubun ortalama aralığı değildir. Her test ayrı iki yönlü 0,05
düzeyinde değerlendirilmiştir; çoklu test düzeltmesi yapılmamıştır.
Ara hesaplarda yuvarlanmamış değerler kullanılmıştır.\par}
\end{table}

\Needspace{13\baselineskip}
\begin{outputguidebox}
Tablo~\ref{tab:b11-dogrulanmis-test} içindeki Welch sonucu,
F grubunun yıl sonu ortalamasının M grubundan yaklaşık 0,948 puan
düşük olduğunu gösterir. Önceden seçilen iki yönlü 0,05 düzeyinde
eşit ortalamalar hipotezi reddedilir. Bu ilişki okul, önceki başarı
ve diğer karıştırıcılar için düzeltilmemiştir; biyolojik veya
nedensel bir açıklama değildir.

Eşleştirilmiş sonuçta yıl sonu eksi ilk dönem ortalama farkı yaklaşık
$-0{,}494$ puandır. Sıfır ortalama fark hipotezi aynı anlamlılık
düzeyinde reddedilir. Tablo~\ref{tab:b11-dogrulanmis-betimsel}
içindeki fark standart sapması kullanıldığında
$d_z=\overline{D}/s_D=-0{,}178706$ bulunur. Bu etki, bağımsız
grupların standartlaştırılmış farkı değildir. Dönem notları ve yıl
sonu notu eşdeğer sınav tekrarları olarak varsayılmaz; sonuç bir
öğretim müdahalesinin etkisi diye yorumlanmaz.

Fark dağılımındaki negatif kuyruk ve aykırı değerler nedeniyle
varsayım değerlendirmesi yalnız $p$ değerine dayandırılamaz.
Büyük örneklem, farklı öğrencilerin bağımsızlığı veya iki okulla
sınırlı kaynağın temsil gücü sorununu çözmez. Yazılımların sayısal
uyumu bu varsayımların doğrulandığı anlamına gelmez.
\end{outputguidebox}

\Needspace{10\baselineskip}
\textbf{Raporlama örneği (doğrulanmış çıktılarla).}
``F grubunda 208, M grubunda 187 öğrenci bulunmaktadır. Welch testinde
F eksi M yıl sonu not ortalaması farkı $-0{,}948092$ puandır
($t=-2{,}065057$, serbestlik derecesi 390,569356,
$p=0{,}039577$; \%95 güven aralığı
$[-1{,}850732;\ -0{,}045452]$). Ayrı eşleştirilmiş analizde
395 öğrencinin yıl sonu eksi ilk dönem ortalama farkı $-0{,}493671$
puan bulunmuştur ($t(394)=-3{,}551703$, $p=0{,}000429$;
\%95 güven aralığı $[-0{,}766937;\ -0{,}220405]$,
$d_z=-0{,}178706$). Sıfır notlu 38 kayıt korunmuştur.
Fark grafikleri negatif kuyruk ve aykırı değerler göstermiştir;
normallik doğrulanmış kabul edilmemiştir. Sonuçlar bağımsızlık
varsayımı altında yorumlanmış; okul ve diğer karıştırıcılar için
düzeltme yapılmamıştır. Bulgular nedensel etki veya daha geniş
bir öğrenci evrenine genellenebilirlik kanıtı olarak sunulmamıştır.''


\Needspace{8\baselineskip}
\subsection{Karşılaştırmalı görev: aynı kayıtlar, iki analiz birimi}
\label{sec:b11-karsilastirmali-gorev}

\textbf{Hedef ve teslim.} Gruplararası fark ile öğrenci içi farkı
ayırt edin. Yeni test çalıştırmadan, mevcut çıktılarla beş maddelik bir
karar ve raporlama kaydı hazırlayın. Önce kendi yanıtınızı yazın;
ardından gerekçeli kontrol yanıtını ve puanlama ölçütlerini okuyun.

\begin{enumerate}
  \item \textbf{Soru ve birim.} Welch ve eşleştirilmiş analiz için
  hedeflenen büyüklüğü, sıfır hipotezini ve analizde kullanılan gözlem
  yapısını ayrı yazın. Aynı 395 kaydı kullanmak neden aynı testi
  gerektirmez? Eşleştirilmiş analizde neden 790 bağımsız gözlem yoktur?
  \item \textbf{Analiz tanımı.} Her iki analiz için fark yönünü, test
  yönünü ve güven düzeyini belirtin. Grup ya da ölçüm sırası ters
  çevrildiğinde fark, $t$, iki yönlü $p$ ve $[L;U]$ güven aralığı
  nasıl değişir? Welch sonucu üzerinde sayısal olarak gösterin.
  \item \textbf{Grafik ve veri kararı.} F--M kutu grafiklerinin
  medyanlarını ve örtüşmesini yorumlayın. Grafikten tek başına test
  kararı verilip verilemeyeceğini ve aykırı görünen kayıtları neden
  otomatik olarak silmeyeceğinizi açıklayın.
  \item \textbf{İki kısa rapor.} Her soru için ayrı bir cümlede
  ortalama farkını puan birimiyle, yüzde 95 güven aralığını ve iki
  yönlü $p$ değerini verin. Gruplararası fark ile öğrenci içi farkı
  aynı büyüklük olarak sunmayın.
  \item \textbf{İddia sınırı.} ``Sonuç anlamlı olduğuna göre kaynak
  cinsiyet kategorisi başarıyı belirler; dönem farkı da öğretimin
  etkisidir'' yorumundaki iki çıkarımı ayrı düzeltin. Okul, önceki
  başarı ve diğer karıştırıcılar için ne yapıldığını belirtin.
\end{enumerate}

\Needspace{8\baselineskip}
\textbf{Gerekçeli kontrol yanıtı.}
\begin{enumerate}
  \item Welch'in hedefi $\mu_{G3,F}-\mu_{G3,M}$ ve sıfır hipotezi
  bu farkın sıfır olmasıdır. Gruplar 208 ve 187 farklı öğrencidir;
  öğrenciler arası bağımsızlık varsayılır. Eşleştirilmiş analizde
  $D=G3-G1$, hedef $\mu_D$, sıfır hipotezi $\mu_D=0$'dır.
  Analiz 395 öğrenci farkını kullanır; bir öğrencinin iki notu
  bağımsız iki kayıt sayılmaz. Testi satır sayısı değil, soru ve
  ölçümlerin yapısı belirler.
  \item Yönler F eksi M ve G3 eksi G1; testler iki yönlü, güven
  düzeyi yüzde 95'tir. Sıra tersine dönünce fark ve $t$ işaret
  değiştirir, iki yönlü $p$ değişmez; aralık $[-U;-L]$ olur.
  Welch için yeni fark $0{,}948092$, $t=2{,}065057$ ve aralık
  $[0{,}045452;\ 1{,}850732]$'dir; $p=0{,}039577$ kalır.
  Eşleştirilmiş fark için de aynı dönüşüm geçerlidir.
  \item Medyanlar F için 10, M için 11'dir; dağılımlar belirgin
  biçimde örtüşür. Medyan çizgileri ortalamaları göstermez ve kutu
  grafiği ortalama farkının testinin yerine geçmez. Aykırı işareti
  veri hatasını kanıtlamaz; kayıtlar incelenir, yalnız işaretlendiği
  için silinmez. Ana analizde sıfır notlar korunmuştur.
  \item ``F eksi M yıl sonu ortalama farkı $-0{,}948092$ puandır;
  yüzde 95 güven aralığı $[-1{,}850732;\ -0{,}045452]$ ve
  iki yönlü $p=0{,}039577$'dir.''
  ``Öğrenci içi G3 eksi G1 ortalama farkı $-0{,}493671$ puandır;
  yüzde 95 güven aralığı $[-0{,}766937;\ -0{,}220405]$ ve
  iki yönlü $p=0{,}000429$'dur.''
  Bu cümleler beşinci maddedeki sınırlılıklarla birlikte okunur.
  \item İlk bulgu düzeltilmemiş bir gruplararası ilişkidir;
  biyolojik veya nedensel belirleme göstermez. İkinci bulgu da
  karşılaştırılabilir bir müdahale tasarımı olmadan öğretim etkisi
  olarak yorumlanamaz. Okul, önceki başarı ve diğer karıştırıcılar
  denetlenmemiştir; notlar eşdeğer sınav tekrarları sayılmamıştır.
\end{enumerate}

\Needspace{8\baselineskip}
\textbf{Değerlendirme ölçütleri (10 puan).} Her görev maddesinde
aşağıdaki iki bileşenin her biri 1 puandır. Doğru ve gerekçeli bileşen
1, eksik veya yanlış bileşen 0 puan alır; eşdeğer doğru ifadeler kabul
edilir. Sayısal sonuçlar altı ondalık basamağa uygun yuvarlanabilir.
\begin{enumerate}
  \item İki hedef ve sıfır hipotezini doğru ayırma (1);
  bağımsız gruplar ile 395 eşleşmiş farkın yapısını açıklama (1).
  \item Fark/test yönleri ve güven düzeyini belirtme (1);
  işaret, $p$ ve aralık dönüşümünü doğru gösterme (1).
  \item Medyanları ve örtüşmeyi doğru okuma (1);
  grafiğin test yerine geçmediğini ve otomatik silmenin uygun
  olmadığını açıklama (1).
  \item Welch için doğru fark, birim, aralık ve $p$ (1);
  eşleştirilmiş analiz için aynı dört öğe (1).
  \item Gruplararası nedensel yorumu karıştırıcılarla birlikte
  düzeltme (1); dönem farkını öğretim etkisinden ayırma (1).
\end{enumerate}
Doğru hesaplar nedensel yorum hatasını telafi etmez; böyle bir yanıt
tam puan alamaz. B12'de bu karşılaştırmayı ortalama ile oran ayrımına
taşıyın (sayfa~\pageref{sec:b12-karsilastirmali-gorev}).